\BourbakiStarChapter{致读者}

\begin{enumerate}
\def\labelenumi{\arabic{enumi}.}
\item
  本系列丛书（卷目见第 ix 页和第 x 页）从数学的起点开始，并给出完整的证明。原则上，它不要求读者具备特定的数学知识，而只要求对数学推理有一定的熟悉程度，并具备一定的抽象思维能力。尽管如此，它尤其面向那些至少已掌握大学数学课程第一年或第二年内容的读者。
\item
  我们选择的阐述方法是公理化的和抽象的，并且通常由一般到特殊。这一选择是由本论著的主要目的所决定的，即为整个现代数学体系提供坚实的基础。为此，必须熟悉相当大量的非常一般的概念和原理。此外，证明的要求对主题内容强加了一个严格固定的顺序。因此，除非读者已经具有相当广泛的数学知识，否则某些考虑的用处不会立即显现；否则，他必须有耐心暂缓判断，直到时机成熟。
\item
  为了减轻这一不利之处，我们经常在正文中插入一些例子，这些例子涉及读者可能已经知道但尚未在本系列中讨论过的事实。这样的例子总是放在两个星号之间：*\ldots*。毫无疑问，大多数读者会发现这些例子有助于他们理解正文，并且即使在第一次阅读时也宁愿不略去它们。当然，从纯粹逻辑的观点来看，略去它们不会有任何不利之处。
\item
  本系列分为若干卷（此处称为“书”）。前六本书有编号，并且一般来说，正文中的每一陈述都只假定前面各卷中已经讨论过的结果为已知。这一规则在每本书内部都成立，但为了阐述的方便，这些书不再按连续顺序排列。在每本书（或这些书的每一章）的开头，读者将找到关于其与其他书逻辑关系的精确说明，从而能够确信不存在任何恶性循环。
\item
  每一章的逻辑框架由该章的定义、公理和定理组成。这些是主要需要记住以备后续使用的部分。较不重要的结果以及那些容易从定理推出的结果被标记为“命题”、“引理”、“推论”、“注记”等。那些在第一次阅读时可以略去的内容用小号字印刷。对特别重要的定理的评注偶尔以“附释”的名义出现。
\end{enumerate}

为了避免冗长的重复，有时引入仅在某一章或某一章中的某一节内有效的记号或缩写是方便的（例如，在一章只涉及交换环时，“环”一词总是指“交换环”）。这样的约定总是被明确提及，通常在其出现的章节的开头。

\begin{enumerate}
\def\labelenumi{\arabic{enumi}.}
\setcounter{enumi}{5}
\item
  正文中的某些段落旨在预先警告读者避免严重错误。这些段落在页边用符号 Z（“危险弯道”）标示。
\item
  练习的设计既使读者能够确信自己已消化了正文，又使他注意到那些在正文中没有位置但仍然有趣的结果。最难的练习带有符号 ¶。
\item
  总的来说，我们遵循了普遍接受的术语，除非有充分的理由偏离它。
\item
  我们特别努力始终使用严格正确的语言，而不牺牲简洁性。我们尽可能在正文中提请注意语言的滥用，没有这些滥用，任何数学文本都有变得迂腐、甚至不可读的风险。
\item
  由于原则上正文是对一个理论的教条式阐述，它通常不包含对文献的引用。参考文献汇集在历史注记中，通常位于每章末尾。这些注记在适当之处也包含对理论中未解决问题的指示。
\end{enumerate}

每个历史注记之后的参考书目通常只包含在所述理论的发展中具有最重要意义的书籍和原始论文。它绝不自诩完备；特别是，仅用于确定优先权问题的参考文献几乎总是被省略。

至于练习，我们一般认为没有必要指出它们的来源，因为它们取自许多不同的来源（原始论文、教科书、习题集）。

\par\noindent\textbf{11. 对本系列某一部分的引用如下给出：}\par

\begin{enumerate}
\def\labelenumi{\alph{enumi})}
\item
  如果引用的是同一节中提出的定理、公理或定义，则用其编号引用。
\item
  如果它们出现在同一章的另一节中，则在该引用中也引用该节。
\item
  如果它们出现在同一本书的另一章中，则引用该章和该节。
\item
  如果它们出现在另一本书中，则首先用其标题引用该书。
\end{enumerate}

结果概要用字母 R 引用；因此，集合论，R 表示“集合论的结果概要”。

\BourbakiStarChapter{《数学原理》系列的内容}

\BourbakiSeriesBook{集合论}

\begin{enumerate}
\def\labelenumi{\arabic{enumi}.}
\tightlist
\item
  形式数学的描述。2. 集合论。3. 有序集；基数；自然数。4. 结构。
\end{enumerate}

\BourbakiSeriesBook{代数}

\begin{enumerate}
\def\labelenumi{\arabic{enumi}.}
\tightlist
\item
  代数结构。2. 线性代数。3. 张量代数、外代数、对称代数。4. 多项式与有理分式。5. 域。6. 有序群与有序域。7. 主理想环上的模。8. 半单模与半单环。9. 半双线性型与二次型。
\end{enumerate}

\BourbakiSeriesBook{一般拓扑学}

\begin{enumerate}
\def\labelenumi{\arabic{enumi}.}
\tightlist
\item
  拓扑结构。2. 一致结构。3. 拓扑群。4. 实数。5. 单参数群。6. 实数空间、仿射空间与射影空间。7. 加法群 \(\mathbf{R}^n\) 。8. 复数。9. 实数在一般拓扑学中的应用。10. 函数空间。
\end{enumerate}

\BourbakiSeriesBook{实变函数}

\begin{enumerate}
\def\labelenumi{\arabic{enumi}.}
\tightlist
\item
  导数。2. 原函数与积分。3. 初等函数。4. 微分方程。5. 函数的局部研究。6. 广义泰勒展开。欧拉-麦克劳林求和公式。7. 伽马函数。词典。
\end{enumerate}

\BourbakiSeriesBook{拓扑向量空间}

\begin{enumerate}
\def\labelenumi{\arabic{enumi}.}
\tightlist
\item
  赋值域上的拓扑向量空间。2. 凸集与局部凸空间。3. 连续线性映射的空间。4. 拓扑向量空间中的对偶性。5. 希尔伯特空间：初等理论。词典。
\end{enumerate}

\BourbakiSeriesBook{积分论}

\begin{enumerate}
\def\labelenumi{\arabic{enumi}.}
\tightlist
\item
  凸性不等式。2. 里斯空间。3. 局部紧空间上的测度。4. 测度的扩张。L \(^{p}\) 空间。5. 测度的积分。6. 向量值积分。7. 哈尔测度。8. 卷积与表示。
\end{enumerate}

\BourbakiSeriesTitle{李群与李代数}

\begin{enumerate}
\def\labelenumi{\arabic{enumi}.}
\tightlist
\item
  李代数。2. 自由李代数。3. 李群。4. 考克斯特群与蒂茨系。5. 由反射生成的群。6. 根系。
\end{enumerate}

\BourbakiSeriesTitle{交换代数}

\begin{enumerate}
\def\labelenumi{\arabic{enumi}.}
\tightlist
\item
  平坦模。2. 局部化。3. 分次、滤过与拓扑。4. 相伴素理想与准素分解。5. 整数。6. 赋值。7. 因子。
\end{enumerate}

\BourbakiSeriesTitle{谱理论}

\begin{enumerate}
\def\labelenumi{\arabic{enumi}.}
\tightlist
\item
  赋范代数。2. 局部紧群。
\end{enumerate}

\BourbakiSeriesTitle{微分与解析流形}

结果概要。

\BourbakiStarChapter{引言}

代数学本质上关注的是计算，即在集合的元素上施行“代数运算”，其中最著名的例子由初等算术的“四则运算”给出。

这里不是追溯逐步抽象这一缓慢过程的地方；通过这一过程，代数运算的概念最初局限于自然数和可测的量，逐渐扩大其领域，同时“数”的概念得到了平行的推广，直到超越后者，它被应用于不具有“数值”特征的元素，例如集合的置换（见第一章的历史注记）。毫无疑问，正是这些连续扩展的可能性——其中计算的形式保持不变，而经受这些计算的数学实体的性质却变化很大——导致了现代数学指导原则的逐渐确立，即数学实体本身并不重要；重要的是它们之间的关系（见第一卷）。无论如何，可以肯定的是，代数学远在数学的其他分支之前就达到了这种抽象水平，并且长期以来习惯于被视为对代数运算的研究，而独立于这些运算可以应用于的数学实体。

失去了任何特定特征，通常代数运算所依据的共同概念非常简单：对两个元素 \(a, b\) （属于同一集合 \(\mathbf{E}\) ）施行代数运算，意味着将有序对 \((a, b)\) 与一个明确定义的第三元素 \(c\) （属于集合 \(\mathbf{E}\) ）相关联。换言之，这个概念中除了函数之外别无他物：给定一个代数运算，就是给定一个定义在 \(\mathbf{E} \times \mathbf{E}\) 上且取值于 \(\mathbf{E}\) 的函数；其唯一特殊之处在于，该函数的定义域是两个与函数取值所在集合相同的集合的乘积；对于这样的函数，我们称之为合成律。

除了这些“内律”之外，数学家们（主要是在几何学的影响下）还被引导去考虑另一种类型的“合成律”；这就是“作用律”，其中除了集合 E（它仍然可以说是处于前景位置）之外，还有一个辅助集合 \(\Omega\) 出现，其元素称为算子：该律这一次将有序对 \((\alpha, a)\) ——由一个算子 \(\alpha \in \Omega\) 和一个元素 \(a \in E\) 组成——与 E 的第二个元素 b 相关联。例如，欧几里得空间 E 上给定中心的一个位似，将实数 k（“位似比”，此处为算子）和 E 的一点 A 与 E 的另一点 A' 相关联：这就是 E 上的一个作用律。

按照一般定义（《集合论》第四章第1节第4号），在集合 E 上给定一个或多个合成律或作用律，就在 E 上定义了一个结构；对于这样定义的结构，我们专门保留代数结构这一名称，而对它们的研究就构成了代数学。

代数结构有若干种（《集合论》第四章第1节第4号），它们一方面由定义它们的合成律或作用律来刻画，另一方面由这些律所满足的公理来刻画。当然，这些公理并非任意选取的，而恰恰是应用中出现的绝大多数律所具有的性质，例如结合性、交换性等等。第一章主要致力于阐述这些公理以及由它们推出的一般性结论；此外，还对两种最重要的代数结构类型作了更细致的研究，即群（其中只出现一个合成律）和环（有两个合成律），域结构是环的一种特殊情形。

第一章中还给出了带算子群和带算子环的定义，其中除了合成律之外，还出现一个或多个作用律。最重要的带算子群是模，向量空间是模的一个特例，它们在经典几何和现代分析中都起着重要作用。对模结构的研究源于线性方程组的研究，因此得名线性代数；这方面的结果见第二章。

同样地，最常出现的带算子环称为代数（或超复系）。在第三、四章中，对两类特殊的代数作了详细研究：外代数，它与由之导出的行列式理论一起成为线性代数的一个有力工具；以及多项式代数，它是代数方程理论的基础。

第五章阐述交换域的一般理论及其分类。这一理论的起源是含一个未知数的代数方程的研究；当年产生这一理论的问题如今已几乎只有历史意义，但交换域的理论对代数学仍然是基础性的，它一方面是一代数论的基础，另一方面也是代数几何的基础。

由于自然数集有两个合成律，即加法和乘法，经典算术（或数论）——即对自然数的研究——从属于代数学。然而，与由这两个律所定义的代数结构相关联的，还有由序关系“a 整除 b”所定义的结构；而经典算术的独特之处恰恰在于研究这两个相关联结构之间的关系。这并不是序结构与代数结构通过“整除性”关系相关联的唯一例子：这一关系在多项式环中同样起着重要作用。因此，第六章将对此作一般性研究；这一研究将在第七章中应用于确定某些特别简单的环上的模的结构，特别是“初等因子”理论。

第八、九章致力于更专门的理论，但它们在分析中有许多应用：一方面是半单模与半单环的理论，它与群的线性表示理论密切相关；另一方面是二次型与埃尔米特型的理论，以及与之相伴的群的研究。最后，初等几何（仿射、射影、欧几里得等）在第二、六、九章中从纯代数的观点加以研究：这里不过是用一种新的语言来表达别处已经得到的代数学结果，但这种语言特别适合代数几何和微分几何的后续发展，本章正是作为它们的引论。
